\section{Discussion}\label{sec:discussion}

\paragraph{What is established.}
Reliable copying of a chemical state is here a property of a complete stochastic cycle, proved uniformly over an admitted region, rather than a consequence of deterministic multistability or a favourable fraction of simulated trajectories. That uniformity is what makes iteration legitimate: Lemma~\ref{lem:iteration} needs no independence between siblings or generations, only that every successful daughter is again a state to which the one-cycle theorem applies. Three quite different chemistries meet this standard. The constructed source shows that the standard can be met with fewer than a thousand resident molecules per module when the environment is ideally controlled, and that two modules can be coupled through shared growth and cross-catalysis without losing it. The redundancy theorem shows that local recovery against bounded total incident interaction is the mechanism that keeps the error exponent and the interaction allowance independent of the number of modules, and that for immediate return the logarithmic dependence on the number of modules, the lineage length and the target error is sharp, with the family depth entering linearly. The Semenov result shows that the proof machinery reaches a published eight-species mechanism at a macroscopic volume, through moving rational metrics instead of state enumeration.

\paragraph{The role of external operations.}
None of the three protocols is autonomous. The constructed source relies on an ideal controller for volume, pool activities and post-division reset; the redundancy model prescribes division at a membrane threshold and equal daughter volumes; the Semenov protocol imposes splitting, refilling and a fixed recovery interval. These operations are part of the theorem statements, and the resource costs they entail are stated separately from the resident counts: reservoir capacities of order $0.17$\,pmol per pool and cycle for the constructed source, a growth deadline of order $k/\gamma$ and rate rescaling with $k$ for the redundancy theorem, and a feed quota of $1.2\times10^{19}$ molecules per daughter for the Semenov protocol. Neither result should be read as a claim about vesicle reproduction, selection or evolution; what the results supply to those questions is a precise, checkable notion of copying fidelity and three worked instances of proving it.

\paragraph{Immediate return versus repair.}
The two copying criteria in this paper differ in a way that matters for lower bounds. Immediate return requires admission at birth and is subject to the partition obstruction of Lemma~\ref{lem:partition}, which no chemistry can evade under fair independent allocation. Return after recovery permits the chemistry to repair allocation noise, and the constructed source's $874$-molecule bound and the Semenov result both exploit this in different ways: the former through the certified recovery built into the birth-to-division interval, the latter through an explicit $500$-second recovery. A single-model comparison of the two criteria, with the same chemistry and the same regions, would sharpen the picture and is not carried out here.

\paragraph{Mathematical limits.}
The constants of the redundancy theorem are existence bounds, many orders of magnitude above anything physically meaningful, and the sufficient prefactor degrades as $\gamma\downarrow0$. The constructed source's regions are boxes at a fixed volume and store concentration states, not compositions (Proposition~\ref{prop:ratio}); its theorem covers two modules and does not extend to more without new local certificates, since the controlled-neighbour argument admits one other module. The Semenov certificate holds at nominal parameters and in narrow regions; Section~\ref{sec:sensitivity} shows that its regions do not survive $0.1\%$ changes in several rate groups, so a laboratory statement would require recentred certificates at measured parameters. All three results concern fair complementary allocation; unequal or correlated division laws, unwanted side reactions and shared-resource fluctuations are outside their scope, and the allocation obstruction of Section~\ref{sec:immediate} shows that some division laws destroy copying regardless of the chemistry.

\paragraph{Method.}
The methodological point common to the three results is the separation between proposing a certificate and validating it. Numerical solvers proposed backward vectors on a finite grid and polynomial metrics along a trajectory; exact integer or rational arithmetic then checked the small set of inequalities that the probability theorem actually consumes, and Lean~4 checked the assembly of those inequalities into statements about the literal stochastic laws. A successful simulation, a small residual or a stable Jacobian was never accepted as closure. This division of labour is what allowed a $2.2\times10^7$-row certificate and a $38$-piece polynomial certificate to become theorems with explicit trust boundaries (Appendix~\ref{app:verification}), and it transfers to other mechanisms for which one wants a fidelity guarantee rather than an estimate.
